\section{Three Finite Summation Laws}

The established finite mechanics contains several exact additive or
accumulative structures. Together they support a general architectural
principle:

\begin{quote}
Global relational quantities may arise through lawful sums of finite
local or compositional contributions.
\end{quote}

The examples below act on different mathematical objects. They are
not identified with one another.

\subsection{Incidence accumulation}

Let
\[
x\in\mathbb R^{15}
\]
be a sector coefficient vector.

The transport-induced sector-edge incidence matrix gives
\[
x
\longmapsto
M^{\mathsf T}x
\longmapsto
MM^{\mathsf T}x
=
Qx.
\]

The intermediate vector
\[
M^{\mathsf T}x\in\mathbb R^{30}
\]
assigns to each core edge the sum of the sector coefficients incident
through the transport-induced sector system.

Applying \(M\) returns those accumulated edge values to the
fifteen-sector register.

Thus
\[
Q=MM^{\mathsf T}
\]
is an exact finite local-to-global-to-local incidence sum.

Because
\[
Q
=
14D_0+9D_1+5D_2+4D_3,
\]
the component at a sector basepoint \(v\) is
\[
(Qx)_v
=
14x_v
+
9\sum_{d(v,u)=1}x_u
+
5\sum_{d(v,u)=2}x_u
+
4\sum_{d(v,u)=3}x_u.
\]

The coefficients
\[
14,\quad9,\quad5,\quad4
\]
are exact sector-overlap multiplicities.

They are not fitted long-range weights.

\subsection{Balanced incidence and common aggregation}

Every row of \(M\) has weight \(14\), and every column has weight
\(7\). Therefore
\[
M\one_{30}
=
14\one_{15}
\]
and
\[
M^{\mathsf T}\one_{15}
=
7\one_{30}.
\]

Consequently,
\[
Q\one_{15}
=
98\one_{15}.
\]

The common aggregate is therefore already determined by the balanced
finite incidence structure:
\[
98=14\cdot7.
\]

This is the common mode isolated by the contrast-concentration
theorem.

\subsection{Binary regional composition}

The twelve-section closure construction contains a binary edge
cochain
\[
\delta
\]
and triangle charge
\[
\omega
\]
satisfying
\[
d\delta=\omega.
\]

For any declared mod-two region \(R\) assembled from native
triangles, interior edges occur twice and cancel. Hence
\[
\bigoplus_{\tau\subset R}\omega(\tau)
=
\bigoplus_{e\subset\partial R}\delta(e).
\]

This is a finite cochain identity.

It expresses a global regional receipt as a sum of boundary
contributions.

No continuous differential geometry is required.

\subsection{Integer regional composition}

The same binary closure class admits certified integer parents.

Let
\[
\eta
\]
denote either admitted parent and define
\[
\Omega=d\eta.
\]

For an oriented finite triangle region \(R\),
\[
\sum_{\tau\subset R}\Omega(\tau)
=
\sum_{e\subset\partial R}\eta(e).
\]

Interior oriented edge contributions cancel.

Reduction modulo two recovers the binary law:
\[
\Omega\bmod2=\omega.
\]

Thus the integer regional sum contains signed additive information
that is absent from the binary reduction.

The binary and integer laws are related by an explicit reduction map,
but they are not the same mathematical object.

\subsection{Registered-history winding}

The certified registered-history construction supplies a third form
of finite accumulation.

Thalean time is
\[
\tauT(\gamma)
=
\operatorname{wind}_{r_1}(\gamma),
\]
where \(r_1\) is the unique \(G_{30}\)-visible primitive
registered-history cycle species.

For composable registered histories,
\[
\tauT(\gamma_2\circ\gamma_1)
=
\tauT(\gamma_2)
+
\tauT(\gamma_1).
\]

Under reversal,
\[
\tauT(\gamma^{-1})
=
-\tauT(\gamma).
\]

This is an additive history coordinate.

It is not an incidence sum on \(G_{15}\), and it is not the regional
integer charge on the twelve-section closure complex.

\subsection{Three domains of summation}

The three exact structures may therefore be summarized as

\[
x
\longmapsto
Qx
\]
on the sector register,

\[
\delta
\longmapsto
d\delta=\omega
\]
on the finite closure complex,

and

\[
\gamma
\longmapsto
\tauT(\gamma)
\]
on registered-history cycles.

Their common feature is lawful finite accumulation.

Their domains, codomains, coefficient systems, and provenance are
different.

No theorem in this paper collapses them into a universal sum.

\subsection{QR long range as accumulated local contribution}

The finite examples motivate the following internal architectural
principle:

\begin{quote}
QR long-range structure is sought as the accumulated result of
admissible finite contributions, not as an independently inserted
long-range selector.
\end{quote}

This principle does not imply that every locally admissible
contribution participates in the global sum.

The summation domain must itself be selected by the relevant
compatibility law.

Schematically,
\[
\text{local contribution}
\longrightarrow
\text{global compatibility}
\longrightarrow
\text{admitted summation domain}
\longrightarrow
\text{registered aggregate}.
\]

The unresolved inter-face transport problem discussed in the next
section lies precisely at the global-compatibility step.

\subsection{Boundary}

The exact theorem-side content of this section is limited to the
three established finite sums.

The broader claim that these structures are manifestations of one
physical long-range law is not made.

In particular,
\[
Qx,
\qquad
\sum\Omega,
\qquad
\tauT
\]
must remain distinct unless an explicit crosswalk is proved.
